\subsection{Relaciones con la planitud.}

\begin{proposition}[3.3.1]
\EGAReaderTarget{ega4.c04.prop.3-3-1}{IV.2.3.3.1-fr}
Sean $f:X\to Y$ un morfismo, $\mathcal F$ un $\mathcal O_X$-módulo
cuasicoherente y plano respecto de $f$, y $\mathcal E$ un
$\mathcal O_Y$-módulo cuasicoherente. Si, para todo $y\in Y$, se pone
$\mathcal F_y=\mathcal F\otimes_{\mathcal O_Y}\boldsymbol{k}(y)$, se tiene
\end{proposition}

\begin{env}[3.3.1.1]
\label{IV.2.3.3.1.1-fr}
\[
\operatorname{Ass}(\mathcal E\otimes_{\mathcal O_Y}\mathcal F)
\supset\bigcup_{y\in\operatorname{Ass}(\mathcal E)}
\operatorname{Ass}(\mathcal F_y)
\]
\end{env}

\emph{y los dos miembros son iguales si $Y$ es localmente noetheriano.}

(Por supuesto, $\mathcal F_y$ es un haz sobre la fibra $f^{-1}(y)$, y se
identifica esta fibra con un subespacio de $X$ (\textbf I, 3.6.1).) Como la cuestión es
local sobre $X$ y sobre $Y$, la cuestión se reduce al caso en que $X$ e $Y$ son
afines; la proposición está entonces demostrada en Bourbaki,
\emph{Alg. comm.}, cap.~IV, \S~2, núm.~6, teo.~2.

\begin{corollary}[3.3.2]
\label{IV.2.3.3.2-fr}
Sean $Y$ un preesquema localmente noetheriano sin ciclos primos asociados
inmersos, $f:X\to Y$ un morfismo y $\mathcal F$ un $\mathcal O_X$-módulo
cuasicoherente y plano respecto de $f$. Entonces, para todo
$x\in\operatorname{Ass}(\mathcal F)$, $f(x)$ es un punto maximal de $Y$.
\end{corollary}

Basta aplicar (3.3.1) con $\mathcal E=\mathcal O_Y$, puesto que
$\operatorname{Ass}(\mathcal O_Y)$ es, por hipótesis, el conjunto de los
puntos maximales de $Y$.

